% Tablas con aire: filas más altas para que el contenido no quede apretado.
\renewcommand{\arraystretch}{1.45}

% Caja del teaser de escena: marco fino con el título montado sobre el borde
% superior. El build (build-pdf.sh) convierte el blockquote del teaser en
% \begin{cajateaser}...\end{cajateaser}; en las demás salidas (web, MD) el
% teaser sigue siendo un blockquote estándar.
\usepackage[most]{tcolorbox}
% Título de la caja por idioma; los build-*.sh lo redefinen para otras ediciones.
\providecommand{\cajateasertitulo}{Lo que acabas de ver}
% Argumento #1: el remate («El comentario de esta escena: …»), montado sobre
% el borde inferior derecho — espejo del título del borde superior.
\newtcolorbox{cajateaser}[1]{
  enhanced,
  colback=white, colframe=black!55, boxrule=0.5pt, arc=1.5pt,
  left=10pt, right=10pt, top=11pt, bottom=11pt,
  before skip=1.6em, after skip=1.9em,
  attach boxed title to top left={xshift=7mm, yshift=-\tcboxedtitleheight/2},
  boxed title style={colback=white, frame hidden, boxrule=0pt,
                     left=4pt, right=4pt, top=1pt, bottom=1pt},
  coltitle=black, fonttitle=\scshape\small,
  title={\cajateasertitulo},
  overlay={\node[fill=white, inner xsep=4pt, inner ysep=1pt, anchor=east,
                 font=\small\itshape]
           at ([xshift=-7mm]frame.south east) {#1};}
}

% Portadilla de parte con la introducción en la MISMA página (sin huérfanas).
\newcounter{libroparte}
\newcommand{\partecon}[1]{%
  \clearpage
  \refstepcounter{libroparte}%
  \phantomsection
  \addcontentsline{toc}{part}{Parte \arabic{libroparte}\hspace{1em}#1}%
  \thispagestyle{plain}%
  \vspace*{4.5em}%
  {\centering
   {\Large\scshape Parte \arabic{libroparte}\par}%
   \vspace{1.1em}%
   {\Huge\bfseries #1\par}}%
  \vspace{3.2em}\noindent\ignorespaces
}
